\chapter{Complexity — NARM Compile}
%%% generated by the blueprint pipeline from TCSlib.Complexity.SpaceComplexity.Machines.NARMCompile
\section{Overview}

Declarations of the Lean module \texttt{TCSlib.Complexity.SpaceComplexity.Machines.NARMCompile}.

\section{Declarations}

\begin{theorem}[Agreement of runs under a choice word]
\lean{Complexity.LogProg.ndtm_runWith_eq_runFrom}
\label{Complexity.LogProg.ndtm_runWith_eq_runFrom}
\leanok
\uses{Turing.Cfg, Turing.MultiTapeTM, Turing.MultiTapeTM.runFrom, Turing.MultiTapeTM.runFrom_succ_eq_step, Turing.MultiTapeTM.step, Turing.MultiTapeTM.step_of_halt, Turing.NDTM, Turing.NDTM.runWith, Turing.NDTM.runWith_cons, Turing.NDTM.stepWith, Turing.NDTM.stepWith_of_halt}
\difficulty{5}
Let $N$ be a nondeterministic machine and $M$ a deterministic machine with the same
number of work tapes, symbol alphabet, and state set. Let $X$ be a property of states,
and suppose their transition actions agree for every choice bit and state not satisfying
$X$. If the run of $M$ from a configuration $c$ is in a state not satisfying $X$ at
every time $t<T$, then for every Boolean word $w$ of length at most $T$, running $N$
from $c$ using $w$ gives the same configuration as running $M$ from $c$ for $|w|$ steps.
\end{theorem}

\begin{lemma}[Avoiding states that force a halt]
\lean{Complexity.LogProg.avoid_of_live}
\label{Complexity.LogProg.avoid_of_live}
\leanok
\uses{Turing.Cfg, Turing.MultiTapeTM, Turing.MultiTapeTM.runFrom, Turing.MultiTapeTM.runFrom_add, Turing.MultiTapeTM.runFrom_of_halt, Turing.MultiTapeTM.runFrom_succ_eq_step', Turing.MultiTapeTM.step}
\difficulty{4}
Let $M$ be a multi-tape Turing machine with $K$ work tapes, alphabet $\mathrm{Sym}$, and
state set $\mathrm{St}$, on a fixed input. Let $X$ be a property of states, and suppose
that whenever a configuration is in a state $q$ satisfying $X$, its next step halts. If
the run from configuration $c$ is still live after $T$ steps, then at every time $t<T$,
any state occupied by the machine fails to satisfy $X$.
\end{lemma}

\begin{lemma}[Output extends the initial output]
\lean{Complexity.LogProg.runFrom_output_extends}
\label{Complexity.LogProg.runFrom_output_extends}
\leanok
\uses{Turing.Cfg, Turing.MultiTapeTM, Turing.MultiTapeTM.outputSymbol, Turing.MultiTapeTM.runFrom, Turing.MultiTapeTM.runFrom_succ_eq_step', Turing.MultiTapeTM.step_output}
\difficulty{4}
Let $M$ be a deterministic multi-tape Turing machine, $c$ a configuration on a fixed
input, and $t$ a natural number. The output after $t$ steps is the initial output of $c$
followed by some finite word of symbols.
\end{lemma}

\begin{lemma}[Avoiding states that force a halt]
\lean{Complexity.LogProg.avoid_of_halt}
\label{Complexity.LogProg.avoid_of_halt}
\leanok
\uses{Complexity.LogProg.runFrom_output_extends, Turing.Cfg, Turing.MultiTapeTM, Turing.MultiTapeTM.runFrom, Turing.MultiTapeTM.runFrom_add, Turing.MultiTapeTM.runFrom_of_halt, Turing.MultiTapeTM.runFrom_succ_eq_step', Turing.MultiTapeTM.step}
\difficulty{5}
Let $M$ be a multi-tape Turing machine with $K$ work tapes, alphabet $\mathrm{Sym}$, and
state set $\mathrm{St}$, on a fixed input. Let $X$ be a property of states, and suppose
that whenever a configuration is in a state $q$ satisfying $X$, its next step halts and
appends the symbol $\beta$ to the current output. If a configuration $c$ has, after $T$
steps, output equal to its initial output followed by a symbol $b\ne\beta$, then at
every time $t<T$, any state occupied by the machine fails to satisfy $X$.
\end{lemma}

\begin{lemma}[Space bound from choice-run intervals]
\lean{Complexity.LogProg.spaceUsedWith_le_of_ranges}
\label{Complexity.LogProg.spaceUsedWith_le_of_ranges}
\leanok
\uses{Turing.Cfg, Turing.NDTM, Turing.NDTM.runWith, Turing.NDTM.spaceUsedWith, Turing.NDTM.visitedWith}
\difficulty{4}
Let $N$ be a nondeterministic machine with $K$ work tapes, let $w$ be a choice word, and
let $c$ be an initial configuration. Suppose that for every choice word $u$ of length at
most $|w|$ and every tape $i$, the head position after running with $u$ lies between
integers $\mathrm{lo}(i)$ and $\mathrm{hi}(i)$, inclusive. Then the total number of
distinct work-tape cells visited along the run under $w$ is at most $\sum_i
(\mathrm{hi}(i)-\mathrm{lo}(i)+1)$.
\end{lemma}

\begin{definition}[Choice targets of the compiled program]
\lean{Complexity.LogProg.nchoice}
\label{Complexity.LogProg.nchoice}
\leanok
\uses{Complexity.LogProg.CSt, Complexity.LogProg.NARM, Complexity.LogProg.NIns.targets, Complexity.LogProg.Ph}
Given a nondeterministic register program $A$, the choice-target map sends a compiled
state to no pair of targets unless it is a program state at the start phase whose
instruction is a nondeterministic choice. In that case it returns the pair of successor
labels, ordered as the target for choice $1$ followed by the target for choice $0$.
\end{definition}

\begin{lemma}[Characterization of Choice Targets]
\lean{Complexity.LogProg.nchoice_spec}
\label{Complexity.LogProg.nchoice_spec}
\leanok
\uses{Complexity.LogProg.CSt, Complexity.LogProg.NARM, Complexity.LogProg.NIns.targets, Complexity.LogProg.Ph, Complexity.LogProg.nchoice}
\difficulty{4}
Let $A$ be a nondeterministic register program, and let $q$ be a compiled state. If the
choice-target function at $q$ returns the pair $(l_1,l_0)$, then there is a program
state $l$ such that $q$ is the program state at the initial phase and $A(l)$ is the
choice instruction with targets $l_1$ and $l_0$.
\end{lemma}

\begin{definition}[Nondeterministic program compilation]
\lean{Complexity.LogProg.ncompileTM}
\label{Complexity.LogProg.ncompileTM}
\leanok
\uses{Complexity.LogProg.CSt, Complexity.LogProg.NARM, Complexity.LogProg.NARM.detPart, Complexity.LogProg.Ph, Complexity.LogProg.armProg, Complexity.LogProg.ctr, Complexity.LogProg.goAct, Complexity.LogProg.nchoice, Turing.MultiTapeTM, Turing.NDTM}
Let $A$ be a nondeterministic register program with $m$ registers and $d$ deciders, let
$l_0$ be an initial program label, and let $D$ be a deterministic Turing machine with
$k_D$ work tapes. For each decider, choose a state of $D$ via a map $q_0$. The compiled
nondeterministic Turing machine has $m+k_D$ work tapes and starts in the program state
at $l_0$; whenever $A$ makes a binary choice, its two transitions proceed to the
corresponding successor labels, and otherwise it uses the compiled deterministic
instruction transition, including the simulation of $D$ when required.
\end{definition}

\begin{definition}[Finite compilation of a nondeterministic register program]
\lean{Complexity.LogProg.ncompileFinNDTM}
\label{Complexity.LogProg.ncompileFinNDTM}
\leanok
\uses{Complexity.LogProg.CSt, Complexity.LogProg.NARM, Complexity.LogProg.Ph, Complexity.LogProg.ncompileTM, Turing.FinNDTM, Turing.MultiTapeTM}
Let $m$, $d$, and $k_D$ be nonnegative integers, and let $\Lambda$ be a type of program
labels and $S_D$ a type of states for a deterministic Turing machine $D$ with $k_D$
work tapes over the Boolean alphabet. Suppose the label and state types are finite and
have decidable equality. For a nondeterministic register program $A$ with $m$ registers
and $d$ deciders, an initial label $l_0$, and a choice of initial state $q_0(j)$ for
each decider $j$, the compilation produces a finite nondeterministic Turing machine over
the Boolean alphabet with $m+k_D$ work tapes and state set $CSt(\Lambda \times Ph, SD,
m)$, whose underlying machine is the compiled machine $ncompileTM(A,l_0,D,q_0)$.
\end{definition}

\begin{lemma}[Congruence of compiled transitions]
\lean{Complexity.LogProg.ctr_congr}
\label{Complexity.LogProg.ctr_congr}
\leanok
\uses{Complexity.LogProg.CSt, Complexity.LogProg.RProg, Complexity.LogProg.ctr, Turing.MultiTapeTM}
\difficulty{4}
Let $P$ and $P'$ be register programs with the same call specification at every state,
and let $D$ be a fixed multi-tape machine and $q_0$ a fixed map of initial decider
states. For any compiled control state $q$, if the ordinary transition functions of $P$
and $P'$ agree at every state $s$ for which $q$ is the program-running state at $s$,
then the compiled one-step actions defined by $P$ and $P'$ agree at $q$.
\end{lemma}

\begin{lemma}[Call behavior is independent of choice replacement]
\lean{Complexity.LogProg.armProg_detPart_call}
\label{Complexity.LogProg.armProg_detPart_call}
\leanok
\uses{Complexity.LogProg.NARM, Complexity.LogProg.NARM.detPart, Complexity.LogProg.NARM.detPart_choose, Complexity.LogProg.NARM.detPart_det, Complexity.LogProg.armProg, Complexity.LogProg.insCall}
\difficulty{4}
Let $A$ assign to each label either a deterministic instruction or a binary choice
between two labels, and let $l_0\in\Lambda$ and
$\beta,\beta'\in\{\mathrm{false},\mathrm{true}\}$. The compiled register programs
obtained by replacing every choice instruction with the corresponding return of $\beta$
and of $\beta'$ have identical call maps.
\end{lemma}

\begin{lemma}[Compiled transition without a nondeterministic choice]
\lean{Complexity.LogProg.ncompileTM_tr}
\label{Complexity.LogProg.ncompileTM_tr}
\leanok
\uses{Complexity.LogProg.CSt, Complexity.LogProg.NARM, Complexity.LogProg.NARM.detPart, Complexity.LogProg.NARM.detPart_choose, Complexity.LogProg.NARM.detPart_det, Complexity.LogProg.NIns.targets, Complexity.LogProg.Ph, Complexity.LogProg.armProg, Complexity.LogProg.armProg_detPart_call, Complexity.LogProg.compileTM, Complexity.LogProg.ctr_congr, Complexity.LogProg.insTr, Complexity.LogProg.nchoice, Complexity.LogProg.ncompileTM, Turing.MultiTapeTM}
\difficulty{5}
Let $A$ be a nondeterministic register program, let $l_0$ be its initial label, let $D$
be a deterministic machine with $k_D$ work tapes, and let $q_0$ assign a state of $D$ to
each of the $d$ deciders. For any compiled state $q$ at which $A$ offers no
nondeterministic choice, and for any bits $b$ and $\beta$, the transition of the
compiled nondeterministic machine under $b$ at $q$ equals the transition of the
deterministic compilation obtained by replacing each choice instruction with the return
value $\beta$.
\end{lemma}

\begin{lemma}[Deterministic replacement at a choice state]
\lean{Complexity.LogProg.detPart_step_choose}
\label{Complexity.LogProg.detPart_step_choose}
\leanok
\uses{Complexity.LogProg.CSt, Complexity.LogProg.NARM, Complexity.LogProg.NARM.detPart, Complexity.LogProg.NARM.detPart_choose, Complexity.LogProg.Ph, Complexity.LogProg.armProg, Complexity.LogProg.compileTM, Complexity.LogProg.ctr, Complexity.LogProg.insCall, Complexity.LogProg.insTr, Complexity.LogProg.nchoice, Complexity.LogProg.nchoice_spec, Complexity.LogProg.retAct, Turing.Action.apply, Turing.Cfg, Turing.MultiTapeTM, Turing.MultiTapeTM.step}
\difficulty{4}
Let $A$ be a nondeterministic register program with $m$ registers and $d$ deciders, let
$D$ be a deterministic Turing machine with $k_D$ work tapes, and choose initial states
for its deciders and an answer bit $\beta$. For any input and configuration $c$ of the
compiled machine, if its state is $q$, and the choice map of $A$ is defined at $q$, then
after one machine step the state is halting and the output is the previous output of $c$
followed by $\beta$.
\end{lemma}

\begin{lemma}[Compiled Choice Step]
\lean{Complexity.LogProg.ncompileTM_stepWith_choose}
\label{Complexity.LogProg.ncompileTM_stepWith_choose}
\leanok
\uses{Complexity.LogProg.NARM, Complexity.LogProg.NIns.targets, Complexity.LogProg.aseam, Complexity.LogProg.goAct, Complexity.LogProg.nchoice, Complexity.LogProg.ncompileTM, Complexity.LogProg.seam, Turing.Action.apply, Turing.MultiTapeTM, Turing.NDTM.stepWith}
\difficulty{4}
Let $A$ be a nondeterministic register program, $l_0$ an initial program state, $D$ a
decider machine, and $q_0$ an assignment of initial decider states. Suppose $A(l)$ is
the choice instruction with targets $l_1$ and $l_0'$. For any input word $x$, choice bit
$b$, register valuation $v$, and output word, one step of the compiled machine from the
configuration formed by adjoining blank auxiliary tapes to the arithmetic configuration
at $l$ with valuation $v$ and output word equals the similarly formed configuration at
$l_1$ if $b$ is true and at $l_0'$ if $b$ is false, with the valuation and output
unchanged.
\end{lemma}

\begin{lemma}[Compiled Machine Initial Configuration]
\lean{Complexity.LogProg.ncompileTM_init}
\label{Complexity.LogProg.ncompileTM_init}
\leanok
\uses{Complexity.LogProg.NARM, Complexity.LogProg.NARM.detPart, Complexity.LogProg.armProg, Complexity.LogProg.aseam, Complexity.LogProg.init_eq_aseam, Complexity.LogProg.ncompileTM, Complexity.LogProg.seam, Complexity.LogProg.seam_init, Turing.MultiTapeTM, Turing.NDTM.initCfg}
\difficulty{3}
Let $A$ be a nondeterministic register program, $l_0$ an initial program state, $D$ a
decider machine, $q_0$ an assignment of initial decider states, and $x$ an input word.
The compiled nondeterministic machine's initial configuration on $x$ is the
configuration formed by adjoining blank auxiliary tapes to the initial arithmetic
configuration at $l_0$, with every register equal to zero and the output empty.
\end{lemma}

\begin{definition}[Work-tape head bounds]
\lean{Complexity.LogProg.NBox}
\label{Complexity.LogProg.NBox}
\leanok
\uses{Turing.Cfg}
Let $C$ be a configuration of a machine with $m+k_D$ work tapes, and let $W,B$ be
natural numbers. The head position of each of the first $m$ work tapes lies between $-1$
and $W$, inclusive, and the absolute value of the head position of each of the remaining
$k_D$ tapes is at most $B$.
\end{definition}

\begin{lemma}[Initial Configuration Satisfies the Tape Bounds]
\lean{Complexity.LogProg.nbox_seam}
\label{Complexity.LogProg.nbox_seam}
\leanok
\uses{Complexity.LogProg.CSt, Complexity.LogProg.NBox, Complexity.LogProg.Ph, Complexity.LogProg.aseam, Complexity.LogProg.seam}
\difficulty{1}
Let $m$ and $k_D$ be the numbers of program and auxiliary tapes, respectively, let $x$
be an input word, and let $W,B$ be natural numbers. For any program state $l$, register
valuation $v$, output word, and state set, the configuration obtained by taking the
initial arithmetic configuration with all registers valued according to $v$ and then
adjoining $k_D$ blank auxiliary tapes has each program-tape head position between $-1$
and $W$, and each auxiliary-tape head position has absolute value at most $B$.
\end{lemma}

\begin{theorem}[Live step of a compiled nondeterministic program]
\lean{Complexity.LogProg.ndet_live}
\label{Complexity.LogProg.ndet_live}
\leanok
\uses{Complexity.LogProg.NARM, Complexity.LogProg.NARM.detPart, Complexity.LogProg.NBox, Complexity.LogProg.NPre, Complexity.LogProg.armProg, Complexity.LogProg.arm_step, Complexity.LogProg.aseam, Complexity.LogProg.astep, Complexity.LogProg.avoid_of_live, Complexity.LogProg.compileTM, Complexity.LogProg.compile_correct_from, Complexity.LogProg.detPart_step_choose, Complexity.LogProg.nchoice, Complexity.LogProg.ncompileTM, Complexity.LogProg.ncompileTM_tr, Complexity.LogProg.ndtm_runWith_eq_runFrom, Complexity.LogProg.seam, Turing.MultiTapeTM, Turing.NDTM.runWith}
\difficulty{6}
Let $A$ be a nondeterministic register program, $l_0$ an initial label, $D$ a
deterministic machine with $k_D$ work tapes, $q_0$ a choice of initial decider states,
and $B,W$ natural-number bounds. Suppose the deterministic interpretation of $A$ with
choices replaced by $0$, started at label $l$ with register values $v$ and no recorded
answer, takes one step to label $l'$ and register values $v'$, still with no recorded
answer; suppose the associated precondition holds and $|\mathrm{bits}(v(r)+1)|\le W$ for
every register $r$. Then there is a natural number $S$ such that every run of the
compiled nondeterministic machine from the corresponding initial configuration using a
choice word of length at most $S$ has register-head positions between $-1$ and $W$ and
decider-tape head positions of absolute value at most $B$, and every choice word of
length exactly $S$ reaches the initial configuration associated with $(l',v')$,
preserving the input and output.
\end{theorem}

\begin{theorem}[Halting behavior of a compiled nondeterministic program]
\lean{Complexity.LogProg.ndet_halt}
\label{Complexity.LogProg.ndet_halt}
\leanok
\uses{Complexity.LogProg.Ins, Complexity.LogProg.NARM, Complexity.LogProg.NARM.detPart, Complexity.LogProg.NBox, Complexity.LogProg.NPre, Complexity.LogProg.armProg, Complexity.LogProg.arm_run, Complexity.LogProg.arun, Complexity.LogProg.aseam, Complexity.LogProg.astep, Complexity.LogProg.avoid_of_halt, Complexity.LogProg.compileTM, Complexity.LogProg.compile_correct_from, Complexity.LogProg.detPart_step_choose, Complexity.LogProg.nastep_det, Complexity.LogProg.nchoice, Complexity.LogProg.ncompileTM, Complexity.LogProg.ncompileTM_tr, Complexity.LogProg.ndtm_runWith_eq_runFrom, Complexity.LogProg.pre_detPart, Complexity.LogProg.rrun, Complexity.LogProg.seam, Complexity.LogProg.seam_state_none, Turing.MultiTapeTM, Turing.NDTM.runWith}
\difficulty{6}
Let $A$ be a nondeterministic register program, $l_0$ an initial label, $D$ a
deterministic machine with $k_D$ work tapes, $q_0$ a choice of initial decider states,
and $B,W$ natural-number bounds. Suppose the deterministic interpretation of $A$ with
choices replaced by $0$, started at label $l$ with register values $v$ and no recorded
answer, halts with answer $b_0$; suppose also that the associated precondition holds and
that every register satisfies $|\mathrm{bits}(v(r)+1)|\le W$, where $|\cdot|$ denotes
word length. Then there are a natural number $S$ and a configuration $C_f$ such that
every run of the compiled nondeterministic machine from the corresponding initial
configuration using a choice word of length at most $S$ has register-head positions
between $-1$ and $W$ and decider-tape head positions of absolute value at most $B$, and
every choice word of length exactly $S$ reaches $C_f$, which is halted and has output
equal to the initial output followed by $b_0$.
\end{theorem}
